% Part: lambda-calculus
% Chapter: introduction
% Section: term-revisited

\documentclass[../../../include/open-logic-section]{subfiles}

\begin{document}

\olfileid{lam}{syn}{tr}
\olsection{\texorpdfstring{项作为 $\alpha$-等价类}{项作为 alpha-等价类}}

从现在起，我们将在$\alpha$-等价的意义下考虑项。
这意味着当我们写出一个项时，我们指的是它所在的$\alpha$-等价类。
例如，我们写$\lambd[a][\lambd[b][a c]]$表示所有与之$\alpha$-等价的项的集合，如$\lambd[a][\lambd[b][a
    c]]$、$\lambd[b][\lambd[a][b c]]$等等。

此外，虽然在前面的章节中，诸如 $N, Q$ 的字母用来表示项，但从现在开始，我们用它们来表示类，接下来的研究对象将是这些类而非项。
诸如 $x, y$ 的字母则继续表示变元。

我们还用记号 $\rep{M}$ 表示类~$M$ 中的任意元素，
若需多个则用$\rep{M}[0], \rep{M}[1], etc. $。

我们沿用项的记号以简化表述。对类作如下定义：
\begin{defn}
  \begin{enumerate}
  \item $\lambd[x][N]$ 定义为包含 $\lambd[x][\rep{N}]$ 的类。
  \item $PQ$ 定义为包含 $\rep{P}\rep{Q}$ 的类。
  \end{enumerate}
\end{defn}

不难看出它们是良定义的，因为$\alpha$-转换是相容的。

\begin{defn} \ollabel{def:fv}
  一个$\alpha$-等价类~$M$ 的\emph{自由变元}，记作 $FV(M)$，定义为 $FV(\rep{M})$。
\end{defn}

正如\olref[alp]{thm:fv}所示，由于$FV(\rep{M}[0]) =
FV(\rep{M}[1])$，这是良定义的。

我们也为类上的代入沿用相应记号：
\begin{defn} \ollabel{def:sub}
  $M$ 中对 $y$ 的 $R$ \emph{代入}，记作 $\Subst{M}{R}{y}$，定义为$\Subst{\rep{M}}{\rep{R}}{y}$，
  其中 $\rep{M}$ 和 $\rep{R}$ 是使得该代入有定义的任意代表元。
\end{defn}

根据\olref[alp]{cor:sub}，这也是良定义的。

请注意，这一定义极大地简化了推理。例如：
\begin{align}
  \Subst{\lambd[x][x]}{y}{x} & =\ollabel{eq:1}\\
  &= \Subst{\lambd[z][z]}{y}{x} \ollabel{eq:2}\\
  &= \lambd[z][\Subst{z}{y}{x}] \\
  &= \lambd[z][z]
\end{align}

如果我们仍将其视为项上的代入，\olref{eq:1}是无定义的；但如前所述，我们现在将其视为类上的代入，因此\olref{eq:2}能够成立：我们可以将 $\lambd[x][x]$ 替换为 $\lambd[z][z]$，因为它们属于同一个类。

出于同样的原因，从现在起我们将假定所选的代表元总是满足代入所需的条件。
例如，当我们遇到$\Subst{\lambd[x][N]}{R}{y}$时，我们将假定所选的代表元 $\lambd[x][N]$ 选取得满足 $x \neq y$ 且 $x \notin FV(R)$。

把 $\lambd[x][x]$ 称为``类''略显别扭，我们今后将它们称为$\Lambda$-\emph{项}（或在本部分下文简称为``项''），以区别于我们熟悉的$\lambda$-项。

\begin{editorial}
  我们还不能完全告别项：$\Lambda$-项的定义是基于$\lambda$-项的，
  且我们尚未给出在$\Lambda$-项上定义函数的方法，
  这意味着所有此类函数都必须先在$\lambda$-项上定义，
  然后再``投射''到$\Lambda$-项上，正如我们对代入所做的那样。
  不过，我们假定读者能直观理解如何在$\Lambda$-项上定义函数。
\end{editorial}

\end{document}
